% --- MAGIC COMMENTS: CONFIGURAZIONE EDITOR ---
% !TEX root = ../../main.tex

% ==============================================================
% MAGIC COMMENTS (COMMENTI MAGICI) - SPIEGAZIONE E GUIDA ALL'USO
% ==============================================================
%
% --- COSA SONO E A COSA SERVONO ---
% I Magic Comments sono commenti speciali, direttive testuali particolari ignorate dal compilatore TeX, e lette
% esclusivamente dall'editor (VS Code, TeXstudio, ecc.). Il loro scopo è quello di istruire localmente l'IDE su come
% trattare i file in cui sono inseriti, ad esempio per
% configurare l'ambiente di lavoro (impostazioni di lingua, correzione grammaticale, eccetera) o su come gestire la
% compilazione di progetti multifile. Agiscono nello specifico file in cui sono inseriti.
%
% NOTA: per chi ha studiato C, ricordano un po' le direttive al preprocessore (#define, #undef eccetera), in particolare
% #pragma, ma sono concetti sostanzialmente diversi, in quanto il preprocessore modifica il codice prima di compilarlo,
% mentre i Magic Comments impostano l'editor senza toccare il codice in alcun caso.
%
% --- DOVE E COME VANNO SCRITTI, ED IN QUALE QUANTITÀ ---
% VS Code è un IDE moderno, tecnicamente potrebbero essere scritti in qualsiasi punto del file, in qualsiasi quantità,
% tuttavia è altamente sconsigliato per questioni di compatibilità con altri editor, come ad esempio TeXShop (su Mac,
% il primo ad aver introdotto i magic comments), che hanno un limite hardcoded (rigoroso) tipico di 20 righe all'inizio
% del file. Quindi per questioni di Good Practice, si consiglia di relegare i magic comments alle primissime righe del
% file: si può eventualmente aggiungere una riga di introduzione estetica, qualora si riesca a rientrare nelle prime 20
% righe con i magic comments.
% Questo tipicamente è più che sufficiente ad ogni utente, visto che il numero di magic comments necessari dovrebbe
% essere piuttosto basso (1-5 per file).
%
% --- IL LIMITE DEI MAGIC COMMENTS: QUANDO NON USARLI E PREFERIRE settings.json ---
% Affidarsi ai Magic Comments per configurazioni "globali", quindi che si vogliono usare nella maggiore parte dei i file
% in progetti multi-file, come ad esempio tipicamente la lingua da associare alla correzione ortografica e grammaticale,
%è ritenuto Bad Practice. Costringe infatti a duplicare lo stesso codice in decine di file ("sporcandoli"), e rischia di
% creare conflitti invisibili, visto che si vanno a scavalcare le impostazioni globali che perderebbero di fatto ogni
% utilità: se l'utente le modificasse, non vedrebbe alcun effetto nel file (e sarebbe facile dimenticarne il motivo).
% In questi casi è sempre preferibile centralizzare tutto nel settings.json del Workspace, che è fatto per questo esatto
% scopo.
%
% --- UN CASO PARTICOLARE: % !TEX root - APPROFONDIMENTO TECNICO IMPORTANTE ---
%
% 1. SCOPO: Lo scopo di questo magic comment è indicare all'editor quale sia il file principale (o "root") del progetto
% LaTeX considerato, da passare al compilatore, quando l'utente avvia il processo di compilazione (o build),
% trovandosi all'interno di un file secondario (incluso tramite \input o \include).
%
% 2. IN QUALI FILE VA INSERITO?
% - Nel main.tex: È totalmente inutile e ridondante. L'editor sa già che è il file principale grazie alla presenza di
%   \begin{document}. In questo progetto viene comunque messo in apertura del main.tex per coerenza, consistenza, e
%   volontà del project manager (che li lo ha posto come utile indicazione), sebbene sia superfluo.
% - Nei file figli (es. \input o \include): È proprio per questi che nasce, ed essenziale aggiungerlo per non far
%   fallire la compilazione, lanciandola all'interno di codesti.
%
% 3. COME SPECIFICARE IL PERCORSO (PATH) DEL FILE ROOT? ESEMPI DI SINTASSI
% Il percorso verso il file root deve seguire una sintassi rigida, e sono disponibili due modalità principali per
% specificarlo:
%
% - PERCORSI RELATIVI (Fortemente Raccomandati): nei quali il percorso è calcolato a partire dalla directory del file
%   corrente. Garantiscono la portabilità del progetto (es. tramite Git/GitHub), consentendo la collaborazione, in
%   quanto funzionerebbero su qualsiasi macchina, a patto che la struttura delle cartelle rimanga intatta,
%   indipendentemente dal sistema operativo o dall'IDE utilizzato, o dall'utente. Esempi di sintassi per percorsi
%   relativi:
%   * File root nella stessa cartella de file secondario:           % !TEX root = main.tex
%   * File root in una cartella inferiore (caso da evitare):        % !TEX root = <foldername>/main.tex
%   * File root in una cartella di un livello superiore:            % !TEX root = ../main.tex
%   * File root in una cartella di n livelli superiore:             % !TEX root = ../../../main.tex (es: di 3 livelli
%                                                                   % superiore)
%   * Navigazione a "V" (cartella superiore poi inferiore):         % !TEX root = ../<foldername>/main.tex
%
%   NOTA sui percorsi relativi: nel caso di stesso livello o livello inferiore, è formalmente corretto utilizzare anche
%   la sintassi esplicita % !TEX root = ./main.tex e % !TEX root = ./<foldername>/main.tex rispettivamente, dove
%   "./" istruisce esplicitamente il sistema operativo a "partire dalla cartella corrente". Tuttavia, l'omissione
%   del punto iniziale è considerata lo standard più pulito, compatto e universalmente supportato dai parser o
%   analizzatore sintattico dell'IDE, ossia lo strumento con cui identifica e interpreta il codice).
%
% - PERCORSI ASSOLUTI (Da Evitare): Indicano il path dalla radice del disco, specificandone appunto la posizione esatta
%   nella macchina che elabora il codice. Rompono irrevocabilmente la compilazione se il progetto viene clonato su una
%   macchina diversa.
%   * Esempio di sintassi per percorso assoluto:            % !TEX root = C:/Utenti/Username/Documenti/Progetto/main.tex
%
% - ATTENZIONE A COME SCRIVERE IL PERCORSO: STANDARD UNIX E SLASH (La regola del Backslash) LaTeX è un ecosistema di
%   derivazione Unix. A prescindere dal sistema operativo in uso (incluso Windows, che nativamente utilizza il backslash
%   "\" per specificare i percorsi), i path all'interno dell'ecosistema TeX devono TASSATIVAMENTE utilizzare il "forward
%   slash" o "slash" ("/"). L'uso del "backslash" ("\") è riservato all'invocazione delle macro o comandi (es. \textbf).
%   In particolare inserire un "\" in un magic comment causerà un errore di parsing, poiché il motore lo interpreterà
%   come l'inizio di un comando inesistente.
%
% 4. VS CODE POSSIEDE L'AUTO-RILEVAMENTO DEL main.tex: PERCHÉ USARE % !TEX root COMUNQUE?
% L'estensione LaTeX Workshop per VS Code possiede un motore di Auto-Rilevamento (Root Directory Check) che, aprendo
% l'intera cartella (Workspace), scansiona l'albero delle dipendenze per dedurre autonomamente il file root.
% Tuttavia, l'uso esplicito di "% !TEX root" rimane la Best Practice per più motivi:
% A. Prestazioni: Evita all'IDE di dover parsare l'intero albero genealogico del progetto a ogni avvio, riducendo il
%    carico di calcolo.
% B. Sicurezza: Previene la compilazione errata se un file figlio viene aperto singolarmente (fuori dal Workspace) o
%    processato da altri editor.
% C. Documentazione: Rende esplicita l'architettura del progetto, facilitando la comprensione e la navigazione,
%    soprattutto per nuovi collaboratori.
% D. Compatibilità: Garantisce che il progetto sia compilabile anche in IDE diversi privi di funzionalità di
%    auto-rilevamento.
%
% A riguardo la documentazione ufficiale di LaTeX Workshop riporta, a link:
% https://github.com/james-yu/latex-workshop/wiki/compile#:~:text=Root%20directory%20check
% Root directory check LaTeX Workshop iterates through all .tex files in the root folder of the workspace. The first one
% containing \documentclass[...]{...} (the [...] is optional) or \begin{document} depending on the value of
% latex-workshop.latex.rootFile.indicator, and which includes the file in the active editor is set as the root file.
% File inclusion is computed using a regex based static parsing of .tex files looking for include, input, import and
% variants of them. Note we cannot compute file inclusions if user defined macros are used to include files. To avoid
% parsing all .tex files in the workspace, you can narrow the search by specifying
% latex-workshop.latex.search.rootFiles.include and/or latex-workshop.latex.search.rootFiles.exclude.
